\documentclass[11pt, reqno]{article}

\usepackage{amsmath, amsthm, amssymb}
\usepackage{enumitem}
\usepackage{tcolorbox}
\usepackage{hyperref}
\usepackage{tikz}
\usepackage{tikz-cd}
\usepackage{pgfplots}
\pgfplotsset{compat=1.18}
\usetikzlibrary{arrows.meta}
\usepackage{mathrsfs}
\usepackage{fancyhdr}
\usepackage[bottom=0.75in, top=1in, left=0.5in, right=0.5in]{geometry}
\usepackage{array}   % for \newcolumntype macro
\newcolumntype{L}{>{$}l<{$}}

\theoremstyle{plain}
\newtheorem*{theorem}{Theorem}
\newtheorem*{proposition}{Proposition}
\newtheorem{exercise}{Exercise}
\newtheorem*{lemma}{Lemma}
\newtheorem*{corollary}{Corollary}

\theoremstyle{definition}
\newtheorem*{definition}{Definition}
\newtheorem*{example}{Example}

\theoremstyle{remark}
\newtheorem*{remark}{Remark}

\renewcommand{\phi}{\varphi}
\renewcommand{\epsilon}{\varepsilon}
\renewcommand{\emptyset}{\varnothing}

\newcommand{\RR}{\mathbb{R}}
\newcommand{\ZZ}{\mathbb{Z}}
\newcommand{\NN}{\mathbb{N}}
\newcommand{\CC}{\mathbb{C}}
\newcommand{\QQ}{\mathbb{Q}}

\DeclareMathOperator{\ima}{\text{im}}

\begin{document}

\topmargin=-40pt
\rhead{Henry Woodburn}
\lhead{Math 696}
\renewcommand{\headrulewidth}{1pt}
\renewcommand{\headsep}{20pt}
\thispagestyle{fancy}

{\Huge \bfseries \noindent Product Rule Notes}
\bigbreak
{\Huge Howdy!}

Now that we have learned how to take derivatives of polynomial and exponential functions, we 
can learn how to take derivatives of more complicated functions. 

We may have functions such as 
\begin{equation}\label{ex1}
    f(x) = (x^2 + 67)e^x
\end{equation}
or
\begin{equation}\label{ex2}
    y = (3x^2 + x +1)(2x^3 + x)
\end{equation}
which are made up of functions we know how to take derivatives of, but
so far we don't know how to handle. 

In both of these cases, we can apply something called the product rule. 
This applies any time we can write 
\[
    f(x) = g(x)\cdot h(x).
\]

\begin{theorem}
    \[
        f'(x) = g'(x)h(x) + g(x)h'(x).
    \]
\end{theorem}

In Example \ref{ex1}, we let 
\[
    g(x) = x^2 + 67\qquad h(x) = e^x
\]
so that 
\[
    f(x) = g(x)\cdot h(x) = (x^2 + 67)e^x.
\]

Then 
\[
    g'(x) = 2x \qquad h'(x) = e^x,
\]
and together 
\[
    f'(x) = g'(x)h(x) + g(x)h'(x) = 2x e^x + (x^2 + 67)e^x.
\]

\end{document}